\section{Глава~\ref{sec:nora}. \nt{НОРА}}

Устройство \nt{НОРА}-системы, её два режима, связь со зрачком и перевёрнутое U в поведении измерены прямо; то, что усиление переключает режим выбора и работает как обратная температура, --- выводы главы; выгода от подстройки усиления --- расчёт по модели книги; «\nt{НОРА} --- ручка поиска» и «\nt{НОРА} --- прерывание» --- гипотезы.

{\footnotesize
\setlength{\tabcolsep}{3pt}\renewcommand{\arraystretch}{1.15}
\begin{longtable}{>{\raggedright}p{0.5cm}>{\raggedright}p{4.0cm}>{\raggedright}p{5.6cm}>{\raggedright}p{2.3cm}>{\raggedright\arraybackslash}p{1.7cm}}
\toprule
№ & Утверждение & Опыт и что измерено & Источник & Статус \\
\midrule\endfirsthead
\toprule
№ & Утверждение & Опыт и что измерено & Источник & Статус \\
\midrule\endhead
1 & \nt{НОРА}-нейронов у человека несколько десятков тысяч, почти все в одной группе & Серийные срезы ствола мозга человека, окраска на тирозингидроксилазу: $53\,900$ клеток в голубом пятне и ещё $6\,260$ в соседнем подъядре. & \cite{Baker1989LC} & измерено \\
2 & Выходы расходятся почти по всему мозгу, но части группы отчасти обслуживают разные области & Вирусно-генетическое мечение у мышей: \nt{НОРА}-нейроны голубого пятна получают входы из широких областей мозга и шлют выходы в широкие области мозга и спинного мозга. У крыс: группа, идущая к миндалине, усиливает выучивание страха, а отдельная группа, идущая к префронтальной коре, --- его угашение. & \cite{Schwarz2015LC, Uematsu2017LC, Poe2020} & измерено \\
3 & Фоновый режим: ровные импульсы с частотой от долей до нескольких в секунду, меняющейся с бодрствованием & Крысы, запись нейронов голубого пятна в цикле сна и бодрствования: частота наибольшая в бодрствовании, ниже в медленном сне, почти ноль в быстром; изменения частоты опережают смену состояния. & \cite{AstonJonesBloom1981} & измерено \\
4 & Отрывистый режим: короткая пачка на важное для задачи событие, примерно в момент решения & Обезьяны, задача найти редкую цель: все нейроны голубого пятна отвечают пачкой только на цель, через $\sim90\ms$ после неё и за $\sim200\ms$ до ответа рукой; при плохой работе пачки слабее. В задаче с выбором из двух пачка точнее привязана к ответу, чем к стимулу, и есть и при верных, и при ошибочных ответах. & \cite{AstonJones1994vigilance, Clayton2004LC} & измерено \\
5 & Диаметр зрачка --- контрольная точка \nt{НОРА} & Обезьяны: импульсы голубого пятна, вплоть до отдельных импульсов одной клетки, предсказывают расширение зрачка; похожая, но более поздняя связь есть в холмиках и поясной коре. Мыши: колебания зрачка следуют за активностью и норадреналиновых, и холинергических выходов в коре --- зрачок отражает не одну \nt{НОРА}. & \cite{Joshi2016, Reimer2016} & измерено \\
6 & \nt{НОРА} меняет крутизну передаточной характеристики~\eqref{eq:gain} & Бодрствующие обезьяны, слуховая кора, ионофорез норадреналина: фоновая активность нейронов подавляется сильнее, чем ответ на голоса сородичей, --- отношение сигнала к шуму растёт. Мультипликативная сигмоида~\eqref{eq:gain} --- упрощение, взятое из сетевой модели. & \cite{Foote1975NE, ServanSchreiber1990} & измерено \\
7 & Усиление повышает $d'$ только против шума после усилителя~\eqref{eq:gainsnr} (утверждение~\ref{prop:gainnoise}) & Вывод в главе; в сетевой модели усиление улучшает обнаружение сигнала. Опыта, который отделял бы шум до и после усилителя и проверял бы это различие, нет. & вывод в главе; \cite{ServanSchreiber1990} & теория \\
8 & Граница $g\beta=1$ переключает сеть выбора из «размышляет» в «решила»~\eqref{eq:wtagain} (утверждение~\ref{prop:gainwta}, рис.~\ref{fig:pitchfork}) & Линейный анализ двух тормозящих друг друга групп; вывод в главе. Прямых измерений этого порога в мозге нет. & вывод в главе & теория \\
9 & Коэффициент усиления --- обратная температура выбора~\eqref{eq:softmax} & При логистическом шуме после усилителя вероятность выбора --- softmax с $T=\sigma/g$; вывод в главе. & вывод в главе & теория \\
10 & \nt{НОРА} задаёт, сколько искать & Люди: перед выбором наугад базовый зрачок шире, чем перед выбором лучшего известного; у людей с более широким зрачком склонность искать выше. Крысы: переход в режим случайного выбора управляется входом \nt{НОРА} в переднюю поясную кору. Направление: больше тонической \nt{НОРА} --- больше поиска, то есть меньше действующее $g$; отождествление \nt{НОРА} с $g$ держится на различии фонового и отрывистого режимов. & \cite{JepmaNieuwenhuis2011, Tervo2014stochastic, AstonJones2005} & гипотеза \\
11 & Награда зависит от $g$ как перевёрнутое U; лучшее $g$ меньше в быстро меняющемся мире (утверждение~\ref{prop:explore}, рис.~\ref{fig:invertedU}) & \texttt{models/nora\_explore.py}, $60$ прогонов по $4000$ попыток. Перемена раз в $1000$ попыток: лучшее $g=16$, доля $0.976$; раз в $20$: $g=8$, доля $0.886$; при $g=0.5$ --- $0.77$. Пересчёт совпал с главой. & \texttt{models/\allowbreak nora\_\allowbreak explore.py} & модель \\
12 & Перевёрнутое U в поведении: лучше всего при умеренной фоновой активности с чёткими пачками & Обезьяны, та же задача, что в строке~4: в периоды хорошей работы фоновая частота умеренная, а пачки на цель чёткие; при высокой фоновой частоте пачек нет и ложных ответов больше. Изменения фоновой и вызванной активности следуют за колебаниями качества работы. & \cite{Usher1999LC, AstonJones2005} & измерено \\
13 & \nt{НОРА} сообщает неожиданную неопределённость, \nt{АЦЕТ} --- ожидаемую & Теория байесовского вывода с двумя видами неопределённости; прямого опыта, разделяющего эти сигналы по медиаторам, в цитируемых работах нет. & \cite{YuDayan2005} & гипотеза \\
14 & После резкой перемены люди учатся быстрее, и зрачок расширяется & Люди, задача предсказания с внезапными переменами: короткие изменения зрачка следуют за оценкой вероятности перемены и вместе с базовым зрачком предсказывают, насколько сильно новые данные меняют оценку (скорость обучения); изменение зрачка, не связанное со светом и задачей, меняет эту скорость. Что зрачок здесь показывает именно \nt{НОРА}, --- вывод по косвенным данным строки~5. & \cite{Nassar2012} & измерено \\
15 & Отрывистая пачка работает как прерывание: сеть сбрасывает состояние & Прямого опыта, где пачка \nt{НОРА} сбрасывает состояние сети, нет; измерены только пачки на важные события (строка~4). & \cite{BouretSara2005} & гипотеза \\
16 & Подстройка $g$ или скорости обучения по неожиданности почти не лучше лучших постоянных настроек & \texttt{models/nora\_explore.py}, мир с эпохами по $1000$ попыток (перемена раз в $1000$ и раз в $20$). Лучшее постоянное $g=8$: $0.925$; подстройка $g$: до $0.927$; подстройка скорости обучения: до $0.926$; постоянная скорость $0.4$: $0.928$. Пересчёт совпал с главой. & \texttt{models/\allowbreak nora\_\allowbreak explore.py} & модель \\
\bottomrule
\end{longtable}}
